% Consolidated textbook-ch07 from f43cd09215a13d29f88baa670fa48441ec2a44e8; source fragments preserved.

% === BEGIN textbook-ch07-choices.tex blob 99db547f05249c3cea4f9aafd15f1b1dd0d32e20 ===
% 第7批：《大学物理简明教程》第5版，第7章静电场选择题7.1.1—7.1.11。
% 逐页核对原书/PDF第185—186页。题意作研究性转述，解析独立书写。

\subsection{选择题：高斯定理、导体、势能与电容}

\sourceproblem{7.1.1}{185}{185}
\problemtext 判断以下关于高斯定理的推论何者成立：
高斯面上场强处处为零是否意味着面内无电荷；
面内电荷为零是否意味着场强处处为零；
高斯面场强非零是否意味着面内必有电荷；
面内有净电荷是否意味着总电通量非零。
\workings
高斯定理为
\[
\oint_S\bm E\cdot d\bm S=\frac{Q_{\mathrm{enc}}}{\varepsilon_0}.
\]
即使面上每点$\bm E=0$也只能推出$Q_{\mathrm{enc}}=0$，并不排除内部存在等量异号电荷；
即使面内无电荷仍可有外部电荷产生电场；
高斯面上场强非零也可全部源于外部电荷。
只有面内净电荷非零必然意味着电通量非零，原题选\textbf{D}。
\crosscheck 通量取决于电场\emph{法向分量的面积积分}，并非取决于场强模的积分。

\sourceproblem{7.1.2}{185}{185}
\problemtext 封闭曲面$S$包住点电荷$Q$，将外部点电荷$q$从无穷远移到$S$附近；
判断穿过$S$的电通量和面上各点的场强是否改变。
\workings
外电荷$q$没有穿越$S$，所以$S$所包住的电荷仍为$Q$，
\[
\boxed{\Phi_E=\oint_S\bm E\cdot d\bm S
=\frac{Q}{\varepsilon_0}}
\]
保持不变。
但各点电场是$Q$与外部$q$的场叠加，
$\bm E=\bm E_Q+\bm E_q$，通常随$q$位置而变化，
故选\textbf{D：通量不变、场强分布改变}。
\crosscheck 外部点电荷的总封闭面通量为零，
不等于其在高斯面\emph{每一点}产生的场强为零。

\sourceproblem{7.1.3}{185}{185}
\problemtext 静电平衡的导体内部，判断电场强度、电势和导体表面电势关系。
\workings
若导体内部$\bm E\ne0$，自由载流子就会继续运动，
不可能达到静电平衡，所以
\[
\boxed{\bm E_{\mathrm{inside}}=\bm0,\qquad
\varphi_{\mathrm{inside}}=\varphi_{\mathrm{surface}}}.
\]
内部电势常数，但一般不为零，原题选\textbf{C}。
\crosscheck 电势零点可任意设定；“内部场强零”不能推出“内部电势零”。

\sourceproblem{7.1.4}{185}{185}
\problemtext 半径$R_a<R_b$的两同心均匀带电球面分别带电$Q_a,Q_b$。
在两球面之间距球心$r$处（$R_a<r<R_b$）求电场强度大小。
\workings
以半径$r$的同心球面为高斯面，仅包住内球面电荷$Q_a$。
外球面因球对称性在其内部场强为零，所以
\[
\boxed{E(r)=\frac{|Q_a|}{4\pi\varepsilon_0r^2}},\qquad
R_a<r<R_b.
\]
原题选\textbf{D}（选择项按正电荷方向写为$Q_a/(4\pi\varepsilon_0r^2)$）。
\crosscheck 如果题设讨论代数径向分量则保留$Q_a$的符号，
若讨论\emph{大小}必须取绝对值。

\sourceproblem{7.1.5}{185--186}{185--186}
\problemtext 导电材料或静电场所在区域的体电荷密度增加至原来的2倍，
在相应电场也按比例增加且介电常数不变时，判断静电能量如何变化。
\workings
由静电学的线性性，在几何、边界和介电性质不变时，
$\rho\mapsto2\rho$导致$\bm E\mapsto2\bm E$。
局部能量密度
\[
w_e=\frac12\varepsilon E^2
\quad\Longrightarrow\quad
\boxed{w'_e=4w_e}.
\]
相同体积上的总电场能亦变为4倍，原题选\textbf{C}。
\crosscheck 这里“电荷密度加倍、几何不变”的限定不可丢；若物体尺寸也变化则不能直接断言能量加倍关系。

\sourceproblem{7.1.6}{186}{186}
\problemtext 一半径$R$的带电导体球，其表面自由电荷面密度为$\sigma$，
求球外紧贴表面处的电场强度。
\workings
取跨越导体表面的薄柱形高斯面，
\[
(\bm E_{\mathrm{out}}-\bm E_{\mathrm{in}})\cdot\bm n
=\frac{\sigma}{\varepsilon_0}.
\]
静电平衡时导体内$\bm E_{\mathrm{in}}=0$，因此
\[
\boxed{E(R^+)=\frac{|\sigma|}{\varepsilon_0}}
\]
（带符号的法向分量为$\sigma/\varepsilon_0$），原题选\textbf{A}。
\crosscheck 此题是\emph{导体球}而非均匀带电的非导体实体球；边界条件跨表面存在跃变。

\sourceproblem{7.1.7}{186}{186}
\problemtext 正方形两个对角顶点各放电荷$Q$，另两个对角顶点各放电荷$q$。
若电荷$Q$所受其他三个电荷的合静电力为零，求$Q$和$q$的关系。
\workings
设正方形边长$a$。对一个顶点的$Q$，
来自对角$Q$的库仑力大小
$F_{QQ}=k_eQ^2/(2a^2)$，沿对角线向外。
来自相邻两个$q$的合力沿同一对角线，大小
$F_{Qq}=\sqrt2\,k_e|Qq|/a^2$。
要平衡必须$Q$与$q$异号，再令大小相等，
\[
\frac{Q^2}{2a^2}=\frac{\sqrt2|Qq|}{a^2}
\quad\Longrightarrow\quad
\boxed{Q=-2\sqrt2\,q}.
\]
原题选\textbf{A}。
\crosscheck 对于$Q\ne0$可约去$|Q|$；当所有电荷为零是平凡解，
题意默认非零。关系与正方形边长无关。

\sourceproblem{7.1.8}{186}{186}
\problemtext 半径$R$、总电荷均匀分布在球体内部的实心绝缘球，
判断电场强度大小$E(r)$从球心到无穷远的函数图像。
\workings
高斯定理给
\[
E(r)=
\begin{cases}
\dfrac{|Q|r}{4\pi\varepsilon_0 R^3},&0\le r<R,\\[3pt]
\dfrac{|Q|}{4\pi\varepsilon_0r^2},&r\ge R.
\end{cases}
\]
从球心$E(0)=0$线性增加至表面最大值，再按$r^{-2}$衰减，
因此选\textbf{B}。
\crosscheck 一定分清\emph{均匀带电实体球}与\emph{均匀带电球面}：
后者在球内场强为零，不是本题模型。

\sourceproblem{7.1.9}{186}{186}
\problemtext 判断平行板电容器两极板分别带$+Q$与$-Q$时是否不带电；
公式$C=Q/U$是否意味着$C$随$Q$而变；
保持带电量不变而增大板距时端电压是否减小；
以及固定电容条件下带电量与电压是否成正比。
\workings
净电荷为零不等于两极板不带电；
理想线性电容$C$由几何与介质决定，不因充电量改变；
保持$Q$时增大板距$d$使$C=\varepsilon S/d$减小，
故$U=Q/C$反而增大。
只有
\[
\boxed{Q=CU}
\]
在固定电容条件下表示带电量与电压成正比，
原题选\textbf{D}。
\crosscheck 电容是元件参数，不应从$C=Q/U$误推因果关系。

\sourceproblem{7.1.10}{186}{186}
\problemtext 半径为$R$的中性空心导体球壳，其空腔内距球心$d$处放一点电荷$q$（$d<R$）。
把导体球壳暂时接地再断开接地线，取无穷远电势为零，求球心$O$的电势。
\workings
在接地导体球形腔中，对内部点电荷可用镜像法使球面电势为零：
在相同径向的$r'=R^2/d$处设镜像电荷$q'=-qR/d$。
于球心的电势
\[
\varphi(O)=\frac{1}{4\pi\varepsilon_0}
\left(\frac qd+\frac{q'}{r'}\right)
=\boxed{\frac{q}{4\pi\varepsilon_0}
\left(\frac1d-\frac1R\right)}.
\]
原题选\textbf{D}。
\crosscheck 地线撤去后导体的已感应电荷保持不变，
外部不再引入净自由电荷时，球壳仍取零电势；
镜像电荷位于腔外，不是真实的自由电荷。

\sourceproblem{7.1.11}{186}{186}
\problemtext 带负电的质点电荷$-q$（$q>0$）从两带电平行板之间的电势零位置穿过小孔到达电势$U>0$的板附近，
求静电场对其所做功与动能增量。
\workings
从初态$\varphi_i=0$到末态$\varphi_f=U$，
场力功
\[
W=q_{\mathrm{particle}}(\varphi_i-\varphi_f)
=(-q)(0-U)=\boxed{qU}.
\]
忽略碰撞、阻力或其他非保守力时，
动能定理给$\boxed{\Delta K=qU}$，
原题选\textbf{B}。
\crosscheck 电荷量为负、电势升高，所以电势能减小，
$\Delta U=(-q)U=-qU$，动能增量相反为正。
% === END textbook-ch07-choices.tex ===

% === BEGIN textbook-ch07-blanks.tex blob 1a2e2b3b0cccf772fbf7f205f2e6f31149886625 ===
% 第7批：《大学物理简明教程》第7章静电场填空题7.2.1—7.2.12。
% 原书/PDF第186—187页；逐条独立研究性转述和推导，题图仍待二校。
\subsection{填空题：电场、通量、极化与电势}

\sourceproblem{7.2.1}{186}{186}
\problemtext 描述静电场的两个基本物理量及其一般定义式。
\workings
两个基本物理量为\textbf{电场强度$\bm E$}和\textbf{电势$\varphi$}。
电场强度定义为静止正试探电荷单位电量所受的力，
\[
\boxed{\bm E=\lim_{q_0\to0}\frac{\bm F}{q_0}}.
\]
以无穷远处电势为零的静电场中，
\[
\boxed{\varphi(P)=\frac{W_{P\to\infty}}{q_0}
=-\int_\infty^P\bm E\cdot d\bm l}.
\]
\crosscheck 电势零点在一般问题中可另选；不可把$\varphi=0$作为“无电场”的充分条件。

\sourceproblem{7.2.2}{186}{186}
\problemtext 静电场中若某一区域的电势在空间处处为同一常量，求该区域的电场强度。
\workings
静电场与电势关系为
\[
\boxed{\bm E=-\nabla\varphi}.
\]
若在一个开区域内$\varphi(\bm r)=\varphi_0$，各方向空间偏导数均为零，故
$\boxed{\bm E=\bm0}$。
\crosscheck 此结论是区域内处处电势不变；
单独某一点的电势等于零不能推出该点场强为零。

\sourceproblem{7.2.3}{186}{186}
\problemtext 点电荷$q$置于立方体中心，求穿过立方体任一面的电通量；
随后把该点电荷移到立方体外，求穿过整个立方体封闭表面的总电通量。
\workings
立方体总通量$\Phi=q/\varepsilon_0$。
当电荷位于中心时，由六个面在立方体对称变换下等价，
每面的\emph{有向}通量相同，
\[
\boxed{\Phi_{\mathrm{face}}=\frac{q}{6\varepsilon_0}}.
\]
电荷移至立方体外，封闭面内净电荷为零，故
\[
\boxed{\Phi_{\mathrm{closed}}=0}.
\]
\crosscheck 电荷位于外部后，每个面单独的通量不必为零；
只保证六个面的有向通量总和为零。

\sourceproblem{7.2.4}{186}{186}
\problemtext 绝缘薄橡胶球壳的总电荷$q$均匀分布于球面且保持不变。
球壳膨胀过程中球面越过距球心$r$处的固定观察点，问该点的电场强度大小如何变化。
\workings
在球面越过观察点之前，该点位于带电球壳外，
\[
E_{\mathrm{before}}=\frac{|q|}{4\pi\varepsilon_0r^2}.
\]
球壳半径变得大于$r$后，该点处于均匀带电球壳内部，
由球对称性与高斯定理
$\boxed{E_{\mathrm{after}}=0}$。
故场强由$\boxed{|q|/(4\pi\varepsilon_0r^2)}$变为$\boxed{0}$。
\crosscheck 观察点与球心距离$r$不变；
电荷只随球面均匀稀释，不能误把球面电荷密度降低直接代入外部点场强公式。

\sourceproblem{7.2.5}{186}{186}
\problemtext 一质子沿图7.2.5中$a\to c\to b$路径移动，电场力做功$8.0\times10^{-15}\,\mathrm J$。
求从$b$沿另一条路径返回$a$时电场力的功；
若$\varphi_a=0$，求$\varphi_b$。
\workings
静电场是保守场，电场力的功只取决于两个端点，
因此反向返回的功
\[
\boxed{W_{b\to a}=-8.0\times10^{-15}\,\mathrm J}.
\]
质子带电量$e=1.602176634\times10^{-19}$ C，
由$W_{a\to b}=e(\varphi_a-\varphi_b)$，
\[
\boxed{\varphi_b=\varphi_a-\frac{W_{a\to b}}e
\approx-4.99\times10^4\,\mathrm V}.
\]
\crosscheck 不能用折线或曲线的实际长度计算静电保守力功；
质子正电且场力做正功，终点电势更低。

\sourceproblem{7.2.6}{186--187}{186--187}
\problemtext 同心金属导体球壳，内壳净电荷$+q$、外壳净电荷$-2q$，
整个系统处于静电平衡。求外壳内表面、外表面的电荷量。
\workings
在外金属壳的材料内部取高斯面，
金属静电平衡有$\bm E=0$，故高斯面内净电荷为零：
$q+Q_{\mathrm{outer,in}}=0$，于是
\[
\boxed{Q_{\mathrm{outer,in}}=-q}.
\]
外球壳自身净电荷$-2q$，
\[
\boxed{Q_{\mathrm{outer,out}}=-2q-(-q)=-q}.
\]
\crosscheck 同心分布使面电荷均匀，但求总电荷不需要预先假定面密度均匀。

\sourceproblem{7.2.7}{187}{187}
\problemtext 将一负点电荷自无穷远移近一个起初不带电、\emph{未接地}的孤立金属导体。
求导体内部静电场强与导体电势相对初态的变化。
\workings
在任意静电平衡状态，导体内总电场$\bm E_{\mathrm{in}}=0$，
所以场强\textbf{不变}，仍为零。
但外部负电荷使导体表面自由电荷重新分布，
\emph{整块导体的电势}通常相对无穷远零点降低。
因此
\[
\boxed{\bm E_{\mathrm{in}}=\bm0}.
\]
导体电势相对初态\textbf{减小}。
\crosscheck 电势是全空间各源的叠加，不因导体内场强为零而必等于零；
孤立与接地是不同边界条件。

\sourceproblem{7.2.8}{187}{187}
\problemtext 一导体外侧填满均匀线性各向同性电介质（相对介电常数$\varepsilon_r$），
测得导体表面附近电场强度大小$E$，求表面自由电荷面密度的大小。
\workings
静电平衡时导体内电场零，界面法向位移矢量边界条件为
$\bm D_{\mathrm{out}}\cdot\bm n=\sigma_{\mathrm f}$。
在线性介质内$\bm D=\varepsilon_0\varepsilon_r\bm E$，
所以
\[
\boxed{|\sigma_{\mathrm f}|=\varepsilon_0\varepsilon_r E}.
\]
\crosscheck 若需要带符号的面密度，取$\sigma_f=\bm D\cdot\bm n$；
已知场强\emph{大小}时仅能求电荷面密度大小。

\sourceproblem{7.2.9}{187}{187}
\problemtext 平行板电容器板间完全充满均匀各向同性电介质$\varepsilon_r$，
忽略边缘效应。已知极板自由电荷面密度大小为$\sigma$，
求介质中的$\bm D$与$E$的大小。
\workings
无限大极板近似下边界条件给
\[
\boxed{D=\sigma},\qquad
\boxed{E=\frac{\sigma}{\varepsilon_0\varepsilon_r}}.
\]
方向由正极板指向负极板。
\crosscheck 这里$\sigma$表示\emph{自由}面电荷密度；
介质表面的极化束缚电荷已隐含在$\varepsilon_r$中，
不能再重复添加。

\sourceproblem{7.2.10}{187}{187}
\problemtext 两张无限大均匀带电平行平面位于真空，
面电荷密度分别为$+\sigma$和$+2\sigma$，板间距$d$。
求两平面之间的电场强度大小。
\workings
每一无限平面在任一侧产生大小$\sigma_i/(2\varepsilon_0)$的场，
方向远离正带电平面。
在两板之间两电场反向叠加，
\[
\boxed{E=\left|\frac{2\sigma}{2\varepsilon_0}
-\frac{\sigma}{2\varepsilon_0}\right|
=\frac{\sigma}{2\varepsilon_0}}.
\]
\crosscheck 该值与$d$无关；
不要将两板内场强错误相加为$3\sigma/(2\varepsilon_0)$。

\sourceproblem{7.2.11}{187}{187}
\problemtext 无外电场时分子正负电荷中心重合的电介质称什么；
施加外电场后分子正负电荷中心相对移动形成何种极化。
\workings
正负电荷中心原本重合的分子无永久电偶极矩，
属于\textbf{非极性电介质}。
外电场使电子云与原子核的中心相对错开，
产生感生电偶极矩，称为\textbf{位移极化}。
填空依次为：\textbf{非极性电介质；位移极化}。
\crosscheck 具有永久偶极矩的极性分子在外电场中还可产生转向极化；
“位移极化”不应误写成“转向极化”。

\sourceproblem{7.2.12}{187}{187}
\problemtext 半径$r_1=5\,\mathrm{cm}$、$r_2=20\,\mathrm{cm}$的同心均匀带电球面，
电荷分别$q_1=+3\times10^{-8}$ C、$q_2=-6\times10^{-8}$ C。
取无穷远零电势，求两球面之间电势为零的球面半径。
\workings
当$r_1<r<r_2$，内球面电势为$k_eq_1/r$，
外球面在其内部贡献常数$k_eq_2/r_2$，所以
\[
\varphi(r)=\frac{1}{4\pi\varepsilon_0}
\left(\frac{q_1}{r}+\frac{q_2}{r_2}\right).
\]
令$\varphi=0$得
\[
\boxed{r=-\frac{q_1}{q_2}r_2
=\frac{3}{6}\cdot20\,\mathrm{cm}
=10\,\mathrm{cm}}.
\]
\crosscheck $5<10<20$ cm落在题设区域中；
如果把外壳在壳内的电势也写成$k_eq_2/r$会得到错误答案。
% === END textbook-ch07-blanks.tex ===

% === BEGIN textbook-ch07-solutions-001-008.tex blob 038f4df77e7106cf0b055d34371828bccad9dd0e ===
% 第7批：《大学物理简明教程》第7章解答题7.3.1—7.3.8。
% 原书/PDF第187页，用户扫描底本；独立理论推导、数值二校待继续。
\subsection{解答题：库仑力与高斯定理}

\sourceproblem{7.3.1}{187}{187}
\problemtext 三个相同的非零点电荷$q$分别固定于边长$a$的等边三角形三个顶点。
(1) 另在三角形中心放什么电荷可使各顶点电荷受到的合力为零？
(2) 平衡条件是否依赖边长？
\workings
对任意顶点电荷，另外两个顶点电荷产生的斥力大小都为
$F_0=k_eq^2/a^2$，夹角为$60^\circ$，
合力沿从中心指向该顶点的径向，大小
\[
F_{\mathrm{vertex}}=2F_0\cos30^\circ
=\sqrt3\,\frac{k_eq^2}{a^2}.
\]
三角形中心到任一顶点的距离为$r=a/\sqrt3$。
为抵消向外合力，中心电荷$Q$必须与$q$异号，且
\[
\frac{k_e|qQ|}{(a/\sqrt3)^2}
=\sqrt3\,\frac{k_eq^2}{a^2}.
\]
得到
\[
\boxed{Q=-\frac{q}{\sqrt3}}.
\]
式中边长$a$约去，因此平衡条件与$a$无关。
\crosscheck 中心电荷本身受三个顶点电荷的合电场为零，
也不受净力。这个平衡属于静力平衡，未断言动力学稳定性。

\sourceproblem{7.3.2}{187}{187}
\problemtext 真空中面积为$S$、相距$d$的两块大平行金属板分别带电$+q$与$-q$。
有人用点电荷公式$F=k_eq^2/d^2$计算板间吸引力，
另有人用$F=qE$和$E=q/(\varepsilon_0 S)$得到$F=q^2/(\varepsilon_0 S)$。
两种说法是否正确？求正确板间力。
\workings
点电荷公式不适用于大面积均匀带电板。
而$E=q/(\varepsilon_0 S)$是\emph{两块板共同作用}的板间总电场；
计算某一板所受的外力时，应只取\emph{另一块板}在该处产生的电场。
无限大均匀带电板单独在每侧产生
$E_{\mathrm{other}}=\sigma/(2\varepsilon_0)$，
$\sigma=q/S$，所以
\[
\boxed{F=qE_{\mathrm{other}}
=\frac{q^2}{2\varepsilon_0 S}}.
\]
方向是两板互相吸引。
\crosscheck 在忽略边缘效应且极板带电量固定的模型中，该力不依赖$d$；
也可由平板电场能的虚功导出相同结果。

\sourceproblem{7.3.3}{187}{187}
\problemtext 两个相同小球质量均为$m$、带电量相同为$q$，
由长度均为$l$的绝缘细线悬挂在同一点。
静止时两线之间夹角$2\theta$（球半径可忽略），求电荷量大小。
\workings
两球之间的距离
\[
r=2l\sin\theta.
\]
对一个球受力平衡，设线张力为$T_s$，
\[
T_s\cos\theta=mg,\qquad
T_s\sin\theta=\frac{k_eq^2}{r^2},
\quad k_e=\frac1{4\pi\varepsilon_0}.
\]
两式相除：
\[
mg\tan\theta=\frac{k_eq^2}{4l^2\sin^2\theta},
\]
故
\[
\boxed{|q|=2l\sin\theta
\sqrt{4\pi\varepsilon_0mg\tan\theta}}.
\]
\crosscheck 两球同号才互斥并分开；
$q^2$的量纲应为$\mathrm{C^2}$，开平方得库仑。

\sourceproblem{7.3.4}{187}{187}
\problemtext 一根长$L=15.0\,\mathrm{cm}$的绝缘均匀带电细杆$AB$，
线电荷密度$\lambda=5.0\times10^{-9}\,\mathrm{C/m}$。
(1) 在杆轴线延长线上、距$B$端$d_1=5.0$ cm的$P$点求场强；
(2) 在杆的垂直平分线上、距中点$d_2=5.0$ cm的$Q$点求场强。
\workings
记$k_e=1/(4\pi\varepsilon_0)$。
(1) 沿杆从$B$向$A$量坐标$s\in[0,L]$，
微元到$P$距离$d_1+s$，故
\[
\boxed{E_P=k_e\lambda\int_0^L\frac{ds}{(d_1+s)^2}
=k_e\lambda\left(\frac1{d_1}-\frac1{d_1+L}\right)}.
\]
代入数值得
$\boxed{E_P\approx6.74\times10^2\,\mathrm{N/C}}$，
方向从杆端$B$沿外侧轴向远离正带电杆。
(2) 取杆中点为$x=0$，垂线距离$d_2$，
沿杆方向的场分量因奇偶对称相互抵消。
垂线方向
\[
E_Q=k_e\lambda\int_{-L/2}^{L/2}
\frac{d_2\,dx}{(x^2+d_2^2)^{3/2}}
=\boxed{\frac{k_e\lambda L}{d_2\sqrt{(L/2)^2+d_2^2}}}.
\]
数值为
$\boxed{E_Q\approx1.50\times10^3\,\mathrm{N/C}}$，
沿垂直平分线背离正带电杆。
\crosscheck 两处场强单位均为$\mathrm{N/C}$，
远离杆时均应按总电荷近似为点电荷而衰减。

\sourceproblem{7.3.5}{187}{187}
\problemtext 点电荷$q$位于立方体中心，求穿过一个面的电通量；
若把该点电荷移至立方体一个顶点，求其六个面的通量。
\workings
中心放置时，对称性和高斯定理给
$\boxed{\Phi_{\mathrm{face}}=q/(6\varepsilon_0)}$。
电荷移至某一顶点后，可将八个相同的立方体拼成大立方体，
使电荷位于大立方体中心，故该小立方体向外的总通量为
$\boxed{\Phi_{\mathrm{total}}=q/(8\varepsilon_0)}$
（以顶点处对称极限解释）。
紧贴该顶点的三张面上，库仑电场与其平面相切，
各面的通量为零；
与这三面相对的三张面因对称性等价，
各有
\[
\boxed{\Phi_{\mathrm{opposite}}=\frac{q}{24\varepsilon_0}}.
\]
\crosscheck 三张远侧面合计$3q/(24\varepsilon_0)=q/(8\varepsilon_0)$；
处理顶点处的奇异电荷须说明取对称极限，不能直接对含奇点的封闭面套用无歧义的表面高斯公式。

\sourceproblem{7.3.6}{187}{187}
\problemtext 均匀带电球壳的内、外半径分别为$a=6.0$ cm、
$b=10.0$ cm，带电体密度
$\rho=2.0\times10^{-5}\,\mathrm{C/m^3}$。
求距球心$r=5.0$、$8.0$、$12.0$ cm处的电场强度大小及方向。
\workings
球对称性下
\[
Q_{\mathrm{enc}}(r)=
\begin{cases}
0,&r<a,\\
\dfrac{4\pi\rho}{3}(r^3-a^3),&a\le r\le b,\\
\dfrac{4\pi\rho}{3}(b^3-a^3),&r>b.
\end{cases}
\]
由高斯定理
\[
E(r)=
\begin{cases}
0,&r<a,\\
\dfrac{\rho(r^3-a^3)}{3\varepsilon_0r^2},&a\le r\le b,\\
\dfrac{\rho(b^3-a^3)}{3\varepsilon_0r^2},&r>b.
\end{cases}
\]
因此
\[
\boxed{E(0.050)=0},\quad
\boxed{E(0.080)\approx3.48\times10^4\,\mathrm{N/C}},
\quad
\boxed{E(0.120)\approx4.10\times10^4\,\mathrm{N/C}}.
\]
$\rho>0$，非零电场均沿径向向外。
\crosscheck $r=a$时场强连续为零，$r=b$时内外表达式相同，
说明均匀\emph{体}电荷分布没有面电荷导致的跳变。

\sourceproblem{7.3.7}{187}{187}
\problemtext 两无限长同轴均匀带电圆柱面半径分别$R_1<R_2$，
单位长度所带电荷分别为$+\lambda$、$-\lambda$。
分别求$r<R_1$、$R_1<r<R_2$、$r>R_2$的场强。
\workings
选半径$r$、长$L$的同轴圆柱形高斯面，
端面通量为零，侧面通量为$E(2\pi rL)$。
按高斯面包围电荷不同，
\[
\boxed{\bm E(r)=
\begin{cases}
\bm0,&r<R_1,\\
\dfrac{\lambda}{2\pi\varepsilon_0 r}\,\bm e_r,&R_1<r<R_2,\\
\bm0,&r>R_2.
\end{cases}}
\]
其中$\lambda>0$，在两圆柱之间沿径向由内柱指向外柱。
\crosscheck 这不要求两圆柱一定为导体，只需无限长圆柱面上电荷轴对称均匀分布；
两柱外的总单位长度电荷为零。

\sourceproblem{7.3.8}{187}{187}
\problemtext 两个无限大平行平面均匀带电，
面电荷密度分别为$+\sigma$、$-\sigma$。
求空间各处电场强度，说明板间和两侧区域的区别。
\workings
单独一张均匀无限带电平面在两侧形成大小
$|\sigma|/(2\varepsilon_0)$、方向由正电荷面向外或向负电荷面内的电场。
在两平面外侧两者反向抵消，
\[
\boxed{\bm E_{\mathrm{outside}}=\bm0}.
\]
在两平面之间，两者方向相同且从正平面指向负平面，
\[
\boxed{E_{\mathrm{between}}=\frac{\sigma}{\varepsilon_0}}.
\]
\crosscheck 此为理想无限平面模型，边缘效应忽略；
若极板有限大小则外部电场并非严格为零。
% === END textbook-ch07-solutions-001-008.tex ===

% === BEGIN textbook-ch07-solutions-009-014.tex blob 1a0b621fa25fef0cd4d85804c54c920752db57eb ===
% 第7批：《大学物理简明教程》第7章静电场解答题7.3.9—7.3.14。
% 原书/PDF第188页，模型从用户上传图核对。
\subsection{解答题：静电势、位移功和电容}

\sourceproblem{7.3.9}{188}{188}
\problemtext 均匀线电荷密度$\lambda$分布在长为$L$的细杆上，
求在杆轴延长线、距最近端$a>0$的点$P$的电势，
以无穷远处电势为零。
\workings
取杆最近端到最远端的坐标$s\in[0,L]$，
微元电荷$dq=\lambda\,ds$距$P$为$a+s$。
故
\[
\varphi_P=\int_0^L\frac1{4\pi\varepsilon_0}
\frac{\lambda\,ds}{a+s}
=\boxed{\frac{\lambda}{4\pi\varepsilon_0}
\ln\frac{a+L}{a}}.
\]
\crosscheck 当$L\ll a$时，
$\ln(1+L/a)\approx L/a$，退化为总电荷$\lambda L$
集中于远处的点电荷电势$k_e\lambda L/a$。

\sourceproblem{7.3.10}{188}{188}
\problemtext 四个相同点电荷
$q_1=q_2=q_3=q_4=1.25\times10^{-8}$ C
分别位于正方形四个顶点，每顶点到中心$O$的距离$r=5.0\times10^{-2}$ m。
(1) 求中心电势；(2) 从无穷远把试验电荷$q_0=1.0\times10^{-9}$ C搬到中心，
求静电场所做功。
\workings
电势是标量，四个电荷的贡献直接相加，
\[
\boxed{\varphi_O=\sum_{j=1}^4\frac{k_eq_j}{r}
=\frac{4k_eq}{r}
\approx8.99\times10^3\,\mathrm V}.
\]
试验电荷的电势能从零增至$q_0\varphi_O$，
静电场做功为其负值，
\[
\boxed{W_{\infty\to O}=-q_0\varphi_O
\approx-8.99\times10^{-6}\,\mathrm J}.
\]
\crosscheck 中心电场因空间对称性可以为零，
但中心电势为正且不为零；不能把“场强零”误读为“电势零”。

\sourceproblem{7.3.11}{188}{188}
\problemtext 直线上点$A,B$分别放置$+q,-q$，
$AB=2R$，中点为$O$。点$C$位于$B$右侧距$B$为$R$处。
试验电荷$q_0$从$O$经上方以$O$为起点的半圆弧移动到$C$，
求过程中的静电场力做功。
\workings
取$O$为坐标原点，
$A=-R$、$B=+R$、$C=+2R$。
在$O$点，
\[
\varphi_O=k_e\left(\frac qR-\frac qR\right)=0.
\]
在$C$点，距$A$为$3R$、距$B$为$R$，
\[
\varphi_C=k_e\left(\frac{q}{3R}-\frac qR\right)
=-\frac{2k_eq}{3R}.
\]
因此
\[
\boxed{W_{O\to C}=q_0(\varphi_O-\varphi_C)
=\frac{2k_eqq_0}{3R}
=\frac{qq_0}{6\pi\varepsilon_0R}}.
\]
\crosscheck 静电力是保守力，与所经过半圆的几何形状无关；
只需要$O,C$两端的电势。

\sourceproblem{7.3.12}{188}{188}
\problemtext 如图7.3.12，均匀正线电荷密度$\lambda$分布在同一平面内的一根弯曲绝缘细线上：
中间是上半圆、半径$R$，两端各接一段长度为$R$的共线水平直线段。
求半圆圆心$O$处的场强矢量与电势（无穷远零电势）。
\workings
设半圆圆心为原点、半圆位于$y>0$。
左右直线段分别处于$[-2R,-R]$及$[R,2R]$，
对中心的水平电场贡献大小相等、方向相反，故合场为零。
半圆弧长微元$d\ell=R\,d\theta$，
每段到圆心距离恒定为$R$。
其场强水平分量对称相消，竖直分量
\[
E_y=-\frac{k_e\lambda}{R}
\int_0^\pi\sin\theta\,d\theta
=-\frac{2k_e\lambda}{R}.
\]
所以
\[
\boxed{\bm E_O=-\frac{2k_e\lambda}{R}\bm e_y}.
\]
电势为标量积分；
半圆贡献$k_e\lambda\pi$，两水平段各贡献
$k_e\lambda\int_R^{2R}dx/x=k_e\lambda\ln2$，
故
\[
\boxed{\varphi_O=k_e\lambda(\pi+2\ln2)}.
\]
\crosscheck 仅半圆弧提供非零的总场强，
但两端直线段都提供非零\emph{电势}；
两种量叠加方式不同。

\sourceproblem{7.3.13}{188}{188}
\problemtext 一均匀带电细圆环半径为$R$、总电荷$Q$，
在环轴线上距环中心$O$为$x$的点$P$求电势。
\workings
环上任意电荷微元$dq$到$P$的距离相同，均为
$s=\sqrt{R^2+x^2}$，所以
\[
\varphi_P=\frac1{4\pi\varepsilon_0}
\int\frac{dq}{\sqrt{R^2+x^2}}
=\boxed{\frac{Q}{4\pi\varepsilon_0\sqrt{R^2+x^2}}}.
\]
\crosscheck $x=0$时$\varphi_O=k_eQ/R$；
$|x|\gg R$时退化为点电荷远场电势$k_eQ/|x|$。

\sourceproblem{7.3.14}{188}{188}
\problemtext 两个平行金属板$A,B$的正对面积为$S=200\,\mathrm{cm^2}$，
板间距$d=2\,\mathrm{cm}$。$B$板接地，$A$板带正电
$Q=7.08\times10^{-7}$ C，忽略边缘效应，取地电势为零。求$A$板电势。
\workings
面积$S=0.0200\,\mathrm{m^2}$、间距$d=0.0200$ m，
平行板电容
\[
C=\frac{\varepsilon_0 S}{d}
=\varepsilon_0
=8.85\times10^{-12}\,\mathrm F
\]
（采用题设真空介电常数的三位有效数字）。
$B$板接地$\varphi_B=0$，
\[
\boxed{\varphi_A=U_{AB}=\frac{Q}{C}
\approx8.0\times10^4\,\mathrm V}.
\]
\crosscheck 在$B$接地且$A$带正电时，
$\varphi_A>\varphi_B$；
只有忽略边缘效应的理想平行板近似下才能直接使用$C=\varepsilon_0S/d$。
% === END textbook-ch07-solutions-009-014.tex ===

% === BEGIN textbook-ch07-solutions-015-020.tex blob 7f997dcdc7debae03a965b10a9b05e609ae26720 ===
% 第7批：《大学物理简明教程》第7章静电场解答题7.3.15—7.3.20。
% 原书/PDF第188页，用户扫描件校对；独立推导及限制条件。
\subsection{解答题：导体静电平衡、电介质与场能}

\sourceproblem{7.3.15}{188}{188}
\problemtext 半径分别为$R_1<R_2$的两同心薄金属球壳，起初内球壳净带电$+q$、
外球壳不带电。
(1) 求外壳内、外表面电荷分布以及外壳电势；
(2) 将外壳接地后再断开地线，求此时外壳电荷分布及电势。
\workings
球对称性及静电平衡使内壳电荷$+q$分布在其外表面。
(1) 在外壳金属内取高斯面，包围电荷为零，故外壳内表面总电荷
$\boxed{Q_{2,\mathrm{in}}=-q}$。
因为外壳原净电荷为零，外壳外表面
$\boxed{Q_{2,\mathrm{out}}=+q}$，
均匀分布于同心球面。
外球壳电势（无穷远为零）
\[
\boxed{\varphi_2=\frac{q}{4\pi\varepsilon_0R_2}},
\]
因为在$r\ge R_2$处系统的总电荷为$q$。
(2) 外壳接地时依然必须有内表面电荷$-q$；
地线使自由电荷流出，外壳外表面净电荷变成零。
因此
\[
\boxed{Q_{2,\mathrm{in}}=-q,\quad
Q_{2,\mathrm{out}}=0,\quad
Q_{2,\mathrm{net}}=-q},\qquad
\boxed{\varphi_2=0}.
\]
\crosscheck 接地后系统在最外部的净包围电荷为零，
故外部静电场与电势均为零。
断开接地线后这个静电平衡解保持（无其他扰动）。

\sourceproblem{7.3.16}{188}{188}
\problemtext 半径为$R_1$的带电金属球总电荷$Q$，
外围包有从$r=R_1$到$r=R_2$的均匀线性电介质球壳，
相对介电常数$\varepsilon_r$，球壳外为真空。
求各区域电场强度、电势以及金属球电势，
以无穷远为零电势。
\workings
对任意半径$r>R_1$的球面，电位移矢量的通量只由自由电荷$Q$决定，
\[
D_r(r)=\frac{Q}{4\pi r^2}.
\]
分段得到
\[
\boxed{\bm E(r)=
\begin{cases}
\bm0,&0\le r<R_1,\\
\dfrac{Q}{4\pi\varepsilon_0\varepsilon_r r^2}\bm e_r,
  &R_1<r<R_2,\\
\dfrac{Q}{4\pi\varepsilon_0 r^2}\bm e_r,&r>R_2.
\end{cases}}
\]
使用$\varphi(r)=\int_r^\infty E_r(s)\,ds$，
外真空区
\[
\boxed{\varphi(r)=\frac{Q}{4\pi\varepsilon_0r}},\quad r\ge R_2.
\]
介质内部
\[
\boxed{\varphi(r)=\frac{Q}{4\pi\varepsilon_0}
\left[\frac1{R_2}+
\frac1{\varepsilon_r}\left(\frac1r-\frac1{R_2}\right)\right]},
\quad R_1\le r\le R_2.
\]
金属球内电势处处等于表面值
\[
\boxed{\varphi_{\mathrm{metal}}=
\frac{Q}{4\pi\varepsilon_0}
\left[\frac1{\varepsilon_rR_1}
+\left(1-\frac1{\varepsilon_r}\right)\frac1{R_2}\right]}.
\]
\crosscheck $\varepsilon_r=1$时退化为真空带电金属球；
在$r=R_2$的两侧电势连续，且$D_r$连续，
$E_r$则因材料介电常数改变而出现跳变。

\sourceproblem{7.3.17}{188}{188}
\problemtext 图7.3.17中，一个中性空心导体$A$的同一空腔内有两个相互绝缘的导体$B,C$。
$B$带净正电，$A$和$C$不带净电。
比较三导体静电平衡时的电势$U_A,U_B,U_C$。
\workings
将空腔外导体$A$作为参考电势$U_A$。
$B$带净正电，其电势相对于边界$A$为正；
带正电的$B$在腔内产生电场。
处于相同腔内的中性浮置导体$C$会极化，
其电势在静电平衡时处处为常数，取值由零净电荷条件决定。
可把$C$的电势看成周围固定电势边界的正权平均；
因此在典型的连通、无其他带电体的构型中
\[
\boxed{U_B>U_C>U_A}.
\]
\crosscheck 上述严格比较依赖$B,C$均独立、互不接触、彼此不屏蔽，
且腔内只有$B$携带非零自由净电荷。
若另施加外部电势边界或其他腔内带电体，则需重新求解静电边值问题。

\sourceproblem{7.3.18}{188}{188}
\problemtext 真空中均匀带电实心球体半径$R$、总电荷$q$，
按电场能密度$w_e=\tfrac12\varepsilon_0E^2$计算其全部静电场能量。
\workings
均匀体电荷密度为$\rho=3q/(4\pi R^3)$。
由高斯定理得到
\[
E(r)=
\begin{cases}
\dfrac{qr}{4\pi\varepsilon_0R^3},&0\le r<R,\\
\dfrac{q}{4\pi\varepsilon_0r^2},&r\ge R.
\end{cases}
\]
用球壳体积元$dV=4\pi r^2dr$积分：
\[
W_{\mathrm{in}}=\frac{\varepsilon_0}{2}
\int_0^R E(r)^2\,4\pi r^2dr
=\frac{q^2}{40\pi\varepsilon_0R},
\]
\[
W_{\mathrm{out}}=\frac{\varepsilon_0}{2}
\int_R^\infty E(r)^2\,4\pi r^2dr
=\frac{q^2}{8\pi\varepsilon_0R}.
\]
所以
\[
\boxed{W_e=W_{\mathrm{in}}+W_{\mathrm{out}}
=\frac{3q^2}{20\pi\varepsilon_0R}}.
\]
\crosscheck 均匀带电\emph{实体}球的静电能量是
$3k_eq^2/(5R)$；若所有电荷只分布在球面，则是$k_eq^2/(2R)$，
二者不要混淆。

\sourceproblem{7.3.19}{188}{188}
\problemtext 真空中平行板电容器正对面积
$S=20.0\,\mathrm{cm^2}$，板间距$d=2.00\,\mathrm{mm}$，
两板电势差$U=400\,\mathrm V$。
忽略边缘效应，求电容器储存的静电能量。
\workings
换算$S=2.00\times10^{-3}\,\mathrm{m^2}$、
$d=2.00\times10^{-3}$ m，
\[
C=\frac{\varepsilon_0S}{d}
=\varepsilon_0\approx8.85\times10^{-12}\,\mathrm F.
\]
因此
\[
\boxed{W_e=\frac12CU^2
=\frac12(8.85\times10^{-12})(400)^2
\approx7.08\times10^{-7}\,\mathrm J}.
\]
\crosscheck 两种等价公式
$W_e=\tfrac12 QU=\tfrac{Q^2}{2C}$也成立；
这里电压单位为伏、能量单位为焦耳。

\sourceproblem{7.3.20}{188}{188}
\problemtext 真空中的球形电容器由两个同心金属球面组成，
内、外半径分别$R_1<R_2$，所带电荷分别为$+Q,-Q$。
求该球形电容器的电容和储能。
\workings
在$R_1<r<R_2$区域，
\[
E(r)=\frac{Q}{4\pi\varepsilon_0r^2}.
\]
两板电势差为
\[
U=\int_{R_1}^{R_2}E(r)\,dr
=\frac{Q}{4\pi\varepsilon_0}
\left(\frac1{R_1}-\frac1{R_2}\right).
\]
所以
\[
\boxed{C=\frac Q U
=\frac{4\pi\varepsilon_0R_1R_2}{R_2-R_1}},
\]
\[
\boxed{W_e=\frac{Q^2}{2C}
=\frac{Q^2}{8\pi\varepsilon_0}
\left(\frac1{R_1}-\frac1{R_2}\right)}.
\]
\crosscheck $R_2\to\infty$时，电容趋于孤立导体球电容$4\pi\varepsilon_0R_1$；
当两壳间距缩小时，电容增大。
% === END textbook-ch07-solutions-015-020.tex ===
